% A content fragment, meant to be \input{} from a master document whose
% preamble loads:
%   \usepackage{tikz}
%   \usetikzlibrary{positioning, calc, fit, backgrounds, arrows.meta}

\section{System Diagram}

\subsection{Human Metabolic System Diagram}

The human metabolic system is a multi-input, single-output system: several inputs act simultaneously on a single output, the measured body mass m. Height, age, and sex (h, a, σ) are fixed parameters, whereas activity (τ, MET), food intake (E_in), and sleep duration (s) are controllable inputs, each entering the Energy Balance through its own path. Driving m toward the healthy body mass m_des therefore requires regulating all controllable inputs together; an input left uncontrolled acts as an unregulated disturbance on D that the remaining inputs cannot reliably compensate for, and the output drifts from its desire.

\begin{figure}[htbp]
\centering
\begin{tikzpicture}[
  node distance = 7mm and 18mm,
  block/.style  = {draw, rectangle, rounded corners=2pt, fill=white,
                   minimum width=41mm, minimum height=13mm, align=center, font=\small},
  sum/.style    = {draw, circle, fill=white, minimum size=6mm, inner sep=0pt, font=\small},
  signal/.style = {-Latex, thick},
  branch/.style = {circle, fill, inner sep=1.2pt},
  system/.style = {draw, rounded corners=6pt, fill=yellow!20, inner sep=6mm},
  lbl/.style    = {font=\footnotesize},
]

% ---- Process blocks (relative placement) ----
% Every block: long name on top, its variable below in gray.
\node[block] (rolling) {7-Day Rolling Average\\ \textcolor{gray}{$\bar{m}$}};
\node[block, below=of rolling] (bmr)   {Resting Metabolic Rate\\ \textcolor{gray}{$BMR$}};
\node[block, below=of bmr]     (act)   {Daily Activity Burn\\ \textcolor{gray}{$E_{act}$}};
\node[block, below=of act]     (tef)   {Thermic Effect of Food\\ \textcolor{gray}{$TEF$}};
\node[block, below=of tef]     (sleep) {Sleep Deprivation Effect\\ \textcolor{gray}{$\delta$}};

\node[block, minimum height=24mm, right=12mm of $(act.east)!0.5!(tef.east)$, yshift=-4mm] (balance)
  {Energy Balance\\ \textcolor{gray}{$D = E_{in} - (BMR + E_{act} + TEF)$}};
\node[block, right=12mm of balance] (mass) {Body Mass Change\\ \textcolor{gray}{$m(t) = \bar{m} + D/\rho$}};
\node[block, above=of mass]    (gly)  {Glycogen + Water Swing\\ \textcolor{gray}{$\Delta M_{gly}$}};
\node[sum,   right=10mm of mass] (sum) {$+$};

% ---- System boundary ----
\begin{scope}[on background layer]
  \node[system, fit=(rolling)(sleep)(balance)(mass)(gly)(sum),
        label={[font=\bfseries]above:Human Metabolic System Model}] (system) {};
\end{scope}

% ---- External inputs (all enter from the left) ----
\coordinate[left=14mm of bmr.west]   (in-bmr);
\coordinate[left=14mm of act.west]   (in-act);
\coordinate[left=14mm of tef.west]   (in-tef);
\coordinate[left=14mm of sleep.west] (in-sleep);
\coordinate[below=9mm of in-sleep]   (in-t);

\draw[signal] ($(in-bmr)+(0,3mm)$) -- node[lbl, above, pos=0.35] {$h, a, \sigma$} ([yshift=3mm]bmr.west);
\draw[signal] (in-act)   -- node[lbl, above, pos=0.35] {$\tau, MET$}       (act.west);
\draw[signal] (in-tef)   -- node[lbl, above, pos=0.35] {$E_{in}$}          (tef.west);
\draw[signal] (in-sleep) -- node[lbl, above, pos=0.35] {$s$}               (sleep.west);
\coordinate (t-branch) at ($(in-t)+(6mm,0)$);
\node[branch] at (t-branch) {};
\draw[signal] (in-t)     -- node[lbl, above, pos=0.1]  {$t$} (in-t -| mass.south) -- (mass.south);
\draw[signal] (t-branch) |- ([yshift=-3mm]bmr.west);

% ---- Internal signals ----
\draw[signal] (bmr.east)   -| node[lbl, above, pos=0.25] {$BMR$}   (balance.north);
\draw[signal] ([yshift=-3mm]act.east) -- node[lbl, above] {$E_{act}$} ([yshift=-3mm]act.east -| balance.west);
\draw[signal] (tef.east) -- node[lbl, above] {$TEF$}   (tef.east -| balance.west);
\draw[signal] (sleep.east) -| node[lbl, above, pos=0.25] {$\delta$} ([xshift=8mm]balance.south);
\coordinate (ein-branch) at ($(tef.west)+(-4mm,0)$);
\node[branch] at (ein-branch) {};
\draw[signal] (ein-branch) |- node[lbl, above, pos=0.7] {$E_{in}$} ($(tef.south)!0.5!(sleep.north)$)
              -| ([xshift=-8mm]balance.south);
\draw[signal] (balance.east) -- node[lbl, above] {$D$} (mass.west);

\draw[signal] (mass.east) -- node[lbl, below] {$m(t)$} (sum.west);
\draw[signal] (gly.east)  -| node[lbl, above, pos=0.25] {$\Delta M_{gly}$} (sum.north);

% ---- Output ----
\coordinate[right=28mm of sum] (out);
\draw[signal] (sum.east) -- (out) node[lbl, right, align=left] {$m$ (measured\\ body mass, kg)};

% ---- Feedback: m into the rolling average, m̄ out to every block that reads mass ----
\coordinate (tap) at ($(sum.east)!0.5!(out)$);
\coordinate (fb-high) at ($(system.north)+(0,12mm)$);
\node[branch] at (tap) {};
\draw[signal] (tap) |- node[lbl, above, pos=0.75] {$m$} (fb-high)
              -| ($(in-bmr |- rolling.west)+(-6mm,0)$) -- (rolling.west);

\coordinate (mbar)     at ([xshift=8mm]rolling.east);
\coordinate (mbar-bmr) at ($(mbar |- rolling.south)!0.5!(mbar |- bmr.north)$);
\node[branch] at (mbar) {};
\node[branch] at (mbar-bmr) {};
\draw[signal] (rolling.east) -- (mbar) -| node[lbl, right, pos=0.75] {$\bar{m}$} (gly.north);
\draw[signal] (mbar) -- (mbar-bmr) -| node[lbl, above, pos=0.25] {$\bar{m}$} (bmr.north);
\draw[signal] (mbar-bmr) |- ($(bmr.south)!0.5!(act.north)$) -| (act.north);

\end{tikzpicture}
\caption[Block diagram of the human metabolic system model.]{Block diagram of the human metabolic system model; the shaded region marks the system boundary. External inputs enter from the left: height, age, and sex ($h, a, \sigma$), activity ($\tau$, $MET$), food intake ($E_{in}$), sleep duration ($s$), and elapsed time ($t$). The Resting Metabolic Rate, Daily Activity Burn, Thermic Effect of Food, and Sleep Deprivation Effect blocks feed the Energy Balance block, whose daily energy deficit $D$ drives the Body Mass Change block. The Glycogen + Water Swing $\Delta M_{gly}$ is added to that body mass to give the measured body mass $m$, which is fed back through the 7-Day Rolling Average as $\bar{m}$ to the Resting Metabolic Rate, Daily Activity Burn, and Glycogen + Water Swing blocks.}
\label{fig_systemdiagram}
\end{figure}

\subsection{Inverse Metabolic System as Feedforward}
